\begin{longtable}{@{}L{1.05cm} L{3.25cm} L{3.9cm} L{4.5cm} C{1.35cm}@{}}
\caption{Ba mươi công trình nhóm đã đọc toàn văn, xếp theo họ phương pháp. Cột cuối cho biết giao thức đánh giá của bài đó có đối chiếu trực tiếp được với giao thức của nhóm hay không.}\label{tab:tongquan}\\
\toprule
\textbf{Mã} & \textbf{Công trình} & \textbf{Mô hình} & \textbf{Kết quả và giao thức} & \textbf{So được?} \\
\midrule\endfirsthead
\multicolumn{5}{@{}l}{\footnotesize\itshape Bảng \thetable{} (tiếp theo)}\\
\toprule
\textbf{Mã} & \textbf{Công trình} & \textbf{Mô hình} & \textbf{Kết quả và giao thức} & \textbf{So được?} \\
\midrule\endhead
\midrule\multicolumn{5}{r@{}}{\footnotesize\itshape tiếp trang sau}\\\endfoot
\bottomrule\endlastfoot
\multicolumn{5}{@{}l}{\cellcolor{rowbg}\textbf{Học sâu dò fQRS}}\\[1pt]
p01 & Zhong 2018\newline{\scriptsize Đây là công trình đầu tiên (theo tuyên bố của chính tác giả) dùng mạng nơ-ron tích chập (CNN) để dò phức bộ QRS thai…} & {\scriptsize CNN 1D nông gồm 3 khối tích chập (Conv1 → BatchNorm → ReLU; Conv2 → Dropout → MaxPool → ReLU; Conv3 → Dropout → MaxPool → ReLU) với 64/128/256 bộ lọc, độ dài bộ lọc bằng 3 cho cả ba lớp (tác giả nói đã thử bộ…} & {\scriptsize CNN đề xuất đạt Precision 75,33\%, Recall 80,54\%, F-measure (F1) 77,85\% và Accuracy 77,38\%; so sánh accuracy bốn bộ phân lớp (Bảng 1): KNN 68,76\%, Naive Bayes 55,52\%, SVM 70,65\%, CNN 77,38\%. Bước SQI chỉ nâng accuracy từ 76,52\%…}\newline{\scriptsize\itshape Chia cố định theo bản ghi một lần duy nhất, không phải tách theo bản ghi (leave-one-record-out), không phải tách theo chủ thể, cũng không phải…} & một phần \\[2pt]
p09 & Xu 2026\newline{\scriptsize Bài giải bài toán tách tín hiệu ECG thai nhi từ ECG bụng đơn kênh (không phải dò đỉnh trực tiếp): EEMD phân rã tín…} & {\scriptsize IMFClassifier1D là CNN 1 chiều nhận đầu vào 1 kênh dài 10.000 mẫu (10 giây ở 1000 Hz), gồm 5 khối conv-BatchNorm-ReLU-pooling với kích thước nhân giảm dần [31, 21, 15, 11, 7] và số kênh đặc trưng hình kim tự…} & {\scriptsize Bộ phân lớp IMF đạt AUC trung bình 0,9282 ± 0,0189 qua tách theo chủ thể trên 20 đối tượng và 3 lần chạy (cao nhất r08 0,9840 ± 0,0035 và Case 12 0,9647; thấp nhất r07 0,8550 ± 0,0161 và Case 9 0,8840), còn cấu hình mạng tốt…}\newline{\scriptsize\itshape Tách theo chủ thể chỉ được áp cho bộ phân lớp IMF chứ không cho bước dò đỉnh R đầu-cuối: 20 đối tượng (15 mô phỏng và 5 lâm sàng), mỗi fold giữ 1…} & \xau{không} \\[2pt]
p10 & Esmaeili 2025\newline{\scriptsize Bài dò QRS thai nhi trực tiếp từ ECG bụng mà không khử ECG mẹ, và không coi đây là bài toán tìm vị trí đỉnh mà là phân…} & {\scriptsize CNN 1 chiều nhận đầu vào 4 (kênh) × 1000 (mẫu), tức đoạn 1 giây ở 1 kHz, gồm 5 lớp tích chập (lớp 1 và 2 dùng stride 2), 7 lớp batch normalization, 3 lớp dropout và 3 lớp fully-connected, kích hoạt ReLU sau…} & {\scriptsize Phương án 1 (25 bản ghi, hàm MAPE, 100 epoch, ngưỡng 0,5) cho kết quả cao nhất: Accuracy 98,36 ± 0,07; Sensitivity 98,79 ± 0,13; Specificity 96,85 ± 0,43; Precision 99,09 ± 0,12; F1 96,35 ± 0,16; MSE 0,01 ± 0. Phương án 5 (60…}\newline{\scriptsize\itshape Không phải tách theo chủ thể cũng không phải tách theo bản ghi, mà là chia ngẫu nhiên 80-20\% (hoặc 33-67\% ở phương án 3), và bài tự mâu thuẫn: Bảng…} & \xau{không} \\[2pt]
p12 & Wahbah 2024\newline{\scriptsize Bài báo xây dựng một khung học sâu phát hiện trực tiếp đỉnh R của thai nhi từ tín hiệu ECG bụng 12 kênh, không tách…} & {\scriptsize Kiến trúc gồm 3 lớp BiLSTM, 2 lớp dropout (0,5) xen giữa, 1 lớp fully connected 2 đầu ra, softmax và lớp phân loại, với đầu vào chuỗi 12 đặc trưng tương ứng 12 kênh bụng. Số tham số tổng: bài không nêu; ngoài…} & {\scriptsize Subject-dependent 5-fold (có rò rỉ): accuracy từng fold 92,5 / 93,4 / 93,5 / 93,2 / 92,6, trung bình 93,0\% với F1 = 0,96; sau khi thêm FECGPP đạt accuracy 94,2\%, F1 = 0,97 và độ nhạy tăng từ 81\% lên 85,8\%. Tách theo chủ thể trên…}\newline{\scriptsize\itshape Bài chạy hai giao thức và báo cáo cả hai: (A) subject-dependent với 5-fold cross-validation, là rò rỉ dữ liệu vì cùng một bệnh nhân xuất hiện ở cả…} & \xau{không} \\[2pt]
\multicolumn{5}{@{}l}{\cellcolor{rowbg}\textbf{Tách nguồn / tái tạo dạng sóng}}\\[1pt]
p03 & Mohebbian 2022\newline{\scriptsize Bài coi việc tách ECG thai nhi khỏi ECG bụng mẹ là bài toán chuyển đổi tín hiệu sang tín hiệu, dùng CycleGAN một chiều…} & {\scriptsize CycleGAN 1D có chú ý: bộ sinh gồm MultiHead self-attention (2 đầu, số chiều nhúng bằng kích thước cửa sổ) → chuẩn hóa và cắt ngưỡng [0,8; 1] → nhân phần tử → Conv1D 13 bộ lọc kernel 3, Conv1D 7 bộ lọc kernel…} & {\scriptsize Chất lượng tái tạo tín hiệu trên A\&D FECG (Bảng I, tách theo chủ thể): R-squared trung bình 98\% [CI 95\%: 97–99\%], theo từng đối tượng R-squared = 0,99 / 0,98 / 0,96 / 0,98 / 0,98; ICC = 0,99 / 0,99 / 0,98 / 0,99 / 0,99; WEDD =…}\newline{\scriptsize\itshape Bài dùng ba giao thức có chất lượng rất khác nhau: trên A\&D FECG là tách theo chủ thể thật sự, lặp 5 lần cho 5 đối tượng — giao thức sạch và là giao…} & \tot{có} \\[2pt]
p04 & Basak 2024\newline{\scriptsize Bài cũng dùng CycleGAN một chiều để chuyển ECG bụng mẹ thành ECG thai nhi, nhưng mục tiêu khác hẳn các bài chỉ dò…} & {\scriptsize CycleGAN 1D với bộ sinh kiểu ResNet: cấu hình tốt nhất là ResNet 13 khối gồm 3 lớp Conv1D hạ mẫu kích thước 16, 32, 64 (mỗi lớp kèm BatchNorm + Dense + ReLU) → n khối ResNet (mỗi khối 2 đơn vị, mỗi đơn vị gồm…} & {\scriptsize Chất lượng tái tạo fECG theo hệ số tương quan Pearson (Bảng 1, trung bình 5 fold): ResNet 13 khối + Basic 0,882 (tốt nhất), ResNet 13 khối + Self 0,871, ResNet 9 khối + Basic 0,864, ResNet 9 khối + Self 0,855, Unet 256 + Basic…}\newline{\scriptsize\itshape Kiểm định chéo 5-fold trên dữ liệu gộp của cả hai bộ, mô tả nguyên văn chỉ gồm một câu duy nhất "The datasets were split using the 5-Fold…} & \xau{không} \\[2pt]
p05 & Arash 2023\newline{\scriptsize Bài giải bài toán tách nguồn mù từ một kênh ECG bụng duy nhất, đồng thời tách tín hiệu trộn AECG thành ECG mẹ (MECG)…} & {\scriptsize DPSS (Dual-Path Source Separation) gồm khối khử nhiễu LSCGAN (bộ sinh DP-LSTM và bộ phân biệt ResBlock) ghép với khối tách nguồn dựa trên mặt nạ (Inception và 3 khối DP-LSTM). Số tham số, kích thước mô hình…} & {\scriptsize Trên FECGSYNDB (tách theo chủ thể, cùng phân phối), F1 trung bình đạt 99,03\% (±0,39) và SI-SNR +14,92 dB (±1,02), tốt nhất là chủ thể 8 với F1 99,35\% và SI-SNR +15,94 dB, yếu nhất là chủ thể 2 với F1 98,24\%. Kiểm thử zero-shot…}\newline{\scriptsize\itshape Mô hình được huấn luyện bằng xác thực chéo tách theo chủ thể chỉ trên FECGSYNDB (9 chủ thể tổng hợp để huấn luyện, 1 chủ thể giữ lại kiểm thử, lặp…} & một phần \\[2pt]
p06 & Lin 2025\newline{\scriptsize Bài xây dựng một mạng hình chữ W gồm hai mạng Attention R2U-Net song song, trong đó nhánh phải tái tạo ECG mẹ, nhánh…} & {\scriptsize Attention R2W-Net gồm hai mạng Attention R2U-Net 1D ghép chữ W, dùng khối tích chập hồi quy có dư (RRC) và cổng chú ý, với số kênh các tầng ghi trong bài là "64, 128, 256, 512, 256 và 1024". Số tham số, kích…} & {\scriptsize Về chất lượng tái tạo tín hiệu, FECGSYN đạt MSE 0,096 ± 0,016 và MAE 0,082 ± 0,006 (theo kịch bản: C0 MSE 0,079 / C1 0,083 / C2 0,081 / C3 0,122 / C4 0,114 / C5 0,097), còn ADFECGDB đạt MSE 0,025 ± 0,004, MAE 0,013 ± 0,003 và…}\newline{\scriptsize\itshape Đây là điểm yếu nghiêm trọng nhất của bài: trên FECGSYNDB chỉ có một lần chia duy nhất — dữ liệu của đối tượng thứ 10 làm tập kiểm thử, các đối…} & \xau{không} \\[2pt]
p07 & Amirreza 2025\newline{\scriptsize Bài chuyển tín hiệu một kênh thành ảnh hai chiều bằng phép nhúng trễ (time-delay embedding, dựa trên định lý Takens)…} & {\scriptsize SCTD-Net gồm SCTD-1Net (chỉ Ursa-Net 2D) và SCTD-2Net (U-Net 1D trích MECG kết hợp Ursa-Net 2D trích FECG). Bảng 3 báo cáo đầy đủ: SCTD-1Net có 7,68 triệu tham số, dung lượng 92,78 MB và 3,91 GFLOPs mỗi lần…} & {\scriptsize SCTD-2Net đánh giá zero-shot đạt Se 97,19\% / PPV 98,80\% / F1 97,97\% trên PCDB (80 kênh) và Se 95,17\% / PPV 97,94\% / F1 96,52\% trên ADFECGDB (17 kênh), vượt SCTD-1Net ở cả ba chỉ số, rõ nhất ở độ nhạy với mức cao hơn khoảng 1,10…}\newline{\scriptsize\itshape Đây là giao thức sạch nhất trong ba bài và cũng là giao thức duy nhất có thể đối chiếu công bằng: huấn luyện 100\% trên dữ liệu tổng hợp FECGSYN…} & một phần \\[2pt]
p11 & Orvas 2025\newline{\scriptsize Bài tách tín hiệu ECG thai nhi từ ECG bụng đơn kênh thu bằng điện cực vải khô - loại không dùng gel, thoải mái khi đeo…} & {\scriptsize Complex UNet (CUNet) là một UNet giá trị phức chạy trên phổ đồ STFT phức. Bài không nêu số tham số, không nêu FLOPs, không nêu độ trễ suy luận, không nêu kích thước mô hình, không nêu số lớp m, kích thước…} & {\scriptsize Trên 100 đối tượng mô phỏng (Bảng II), CUNet đạt PRD 26,2 ± 11,6; PCC 95,9 ± 5,7; F-score 99,8 ± 0,2; SE 99,6 ± 0,4; HRerr 0,40 ± 0,81 bpm, so với UNet F-score 99,3 ± 1,6 / PCC 77,4; AttUNet 98,4 ± 10,1 / PCC 77,4; SVD 92,0 ±…}\newline{\scriptsize\itshape Không phải tách theo chủ thể cũng không phải tách theo bản ghi, nhưng cũng không cần, vì tập huấn luyện là mô phỏng còn tập đánh giá là dữ liệu…} & \tot{có} \\[2pt]
\multicolumn{5}{@{}l}{\cellcolor{rowbg}\textbf{Xử lý tín hiệu cổ điển}}\\[1pt]
p16 & Niknazar 2013\newline{\scriptsize Bài giải quyết bài toán tách ECG thai khỏi tín hiệu bụng khi chỉ có một điện cực duy nhất, bằng cách mở rộng vector…} & {\scriptsize Đây không phải mạng nơ-ron mà là bộ lọc Kalman mở rộng dựa trên mô hình động phi tuyến; bài không nêu tổng số tham số một cách tường minh, cũng không nêu độ trễ, thời gian chạy hay kích thước mô hình. Suy…} & {\scriptsize Mọi con số định lượng đều là mức cải thiện SNR/SIR trên dữ liệu tổng hợp chứ không phải độ chính xác phát hiện đỉnh: khi hỗn hợp gần như sạch (công suất nhiễu -30 dB), mức cải thiện SNR tối thiểu của fECG là 40 dB, và ngay cả…}\newline{\scriptsize\itshape Bài không nêu dung sai khớp nhịp vì chưa bao giờ đo độ chính xác phát hiện đỉnh: không có Se, không có PPV, không có F1 và không có MAE trên bất kỳ…} & \xau{không} \\[2pt]
p17 & Castillo 2018\newline{\scriptsize Bài đo nhịp tim thai từ một kênh ECG bụng duy nhất, với điểm độc đáo là không khử ECG mẹ: thay vì tách fECG khỏi hỗn…} & {\scriptsize Đây không phải mô hình học sâu và không có tham số huấn luyện theo nghĩa mạng nơ-ron, mà là chuỗi xử lý tín hiệu kết hợp gom cụm không giám sát, cài đặt bằng MATLAB ở mức chứng minh khái niệm; bài không nêu…} & {\scriptsize Trên ADFECGDB với dung sai 50 ms: toàn bộ 17 tín hiệu (10.839 fQRS, 10.000 TP, 839 FN, 413 FP) đạt Se 92,26\%, PPV 96,03\%, Acc 88,87\%, F1 94,11\%; riêng Nhóm 1 gồm 11 tín hiệu sạch (6.990 fQRS, 175 FN tức 2,50\%, 98 FP tức 1,40\%)…}\newline{\scriptsize\itshape Dung sai khớp nhịp là 50 ms (nguyên văn: mọi ứng viên lệch quá 50 ms so với nhãn tham chiếu đều bị coi là FP), trùng khớp hoàn toàn với dung sai của…} & \tot{có} \\[2pt]
p19 & Eleonora 2020\newline{\scriptsize Bài không đề xuất mô hình mới mà trả lời một câu hỏi kỹ thuật hẹp: khi dùng bộ lọc thích nghi để xóa ECG mẹ khỏi tín…} & {\scriptsize Bộ lọc thích nghi QRD-RLS, không phải mạng nơ-ron và không dùng học sâu, với hệ số quên lambda = 0,999 và số tạp N = 20 cho mỗi kênh, nghĩa là MR có 60 hệ số còn SR có 20 hệ số. Bài không nêu số tham số của…} & {\scriptsize Trên dữ liệu thật (trung vị), độ chính xác phát hiện fQRS của SR so với MR là 0,62 so với 0,85 ở kênh ngang, 0,69 so với 0,92 ở kênh dọc, 0,61 so với 0,89 ở kênh chéo và 0,70 so với 0,86 ở kênh đơn cực, dẫn đến kết luận SR chỉ…}\newline{\scriptsize\itshape Không có chia train/test, không tách theo chủ thể, không tách theo bản ghi và không k-fold vì đây không phải bài học máy; tuy nhiên lambda, N và…} & \xau{không} \\[2pt]
\multicolumn{5}{@{}l}{\cellcolor{rowbg}\textbf{Đối chuẩn, dữ liệu, thử thách}}\\[1pt]
p13 & Behar 2014\newline{\scriptsize Đây là luận án Tiến sĩ 249 trang, tài liệu nền tảng của lĩnh vực fECG không xâm lấn và là nguồn gốc của chuẩn đánh giá…} & {\scriptsize Luận án không dùng mạng học sâu; các mô hình có tham số gồm ESN với bể chứa M = 90 nơ ron, độ thưa ψ = 20\%, bán kính phổ ρ = 0,4, hệ số rò rỉ α = 0,4 và input scaling γ = 1, LMS với N = 20 và μ = 0,1, RLS với…} & {\scriptsize Chương 5: bSQI đơn lẻ tốt nhất với Ac = 0,969 (trên pcaSQI 0,950, sSQI 0,892, kSQI 0,879, qSQI 0,766, pSQI 0,752, basSQI 0,626), kết hợp 6 chỉ số cho Ac kiểm thử 0,985, nhưng trên MIMIC II hoàn toàn độc lập thì tụt còn Ac =…}\newline{\scriptsize\itshape Dung sai khớp nhịp là ±50 ms, ghi rõ tại trang 59 và nhắc lại ở mục 9.2.2 — đây chính là nguồn gốc chuẩn mà nhóm đang dùng — với Se = TP/(TP+FN),…} & một phần \\[2pt]
p14 & Behar 2014\newline{\scriptsize Chương này so sánh công bằng 15 phương pháp tách fECG trên cùng một bộ dữ liệu, cùng bộ tiền xử lý, cùng bộ dò đỉnh và…} & {\scriptsize Không có mạng nơ ron và không có học sâu — toàn bộ là xử lý tín hiệu cổ điển kết hợp tối ưu siêu tham số, nên số tham số huấn luyện được: không áp dụng. Hệ thống chỉ có 6 con số toàn cục (nbC = 20, nbPC = 2,…} & {\scriptsize Trên set-a với fb = 10 Hz, thứ hạng F1 là FUSE-SMOOTH 96,0\% (Se 95,9, PPV 96,0, HRE 19,1, RRE 6,3), FUSE 95,0, TSpca-ICA 93,0, TSc-ICA 92,6, ICA-TSpca-ICA 92,4, TS-ICA 92,0, TSm-ICA 91,2, ICA-TSpca 89,2, TSlp-ICA 88,4, TSpca…}\newline{\scriptsize\itshape Dung sai khớp nhịp là ±50 ms, và hai hệ thống chấm điểm chạy song song không trùng nhau: F1 kèm Se và PPV được tính trên set-a, tức trong mẫu vì…} & một phần \\[2pt]
p15 & Andreotti 2016\newline{\scriptsize Bài này không đề xuất mô hình mới mà xây dựng một bộ khung kiểm thử chuẩn (stress-testing) cho lĩnh vực đo ECG thai…} & {\scriptsize Không có mạng nơ-ron nào được huấn luyện (ngoại trừ AMesn là mạng echo state, nhưng bài không nêu số nơ-ron hay số tham số của nó); đây là 8 thuật toán xử lý tín hiệu cổ điển và bài không nêu số tham số, kích…} & {\scriptsize Thí nghiệm 1 (số kênh cho BSS): với 8 kênh, F1 trung bình đạt 97,46\% (BSSica-JADE), 97,40\% (BSSpca) và 97,22\% (FAST-ICA); khi Nch = 16, bước giảm chiều bằng PCA nâng F1 trung bình của ICA từ 95,3\% lên 97,3\%, và F1 gần như bão…}\newline{\scriptsize\itshape Không có chia huấn luyện/kiểm thử vì cả 8 phương pháp đều là xử lý tín hiệu không giám sát chạy trực tiếp trên từng bản ghi, nên không tồn tại khái…} & một phần \\[2pt]
p20 & Gari 2014\newline{\scriptsize Đây là bài xã luận mở đầu số đặc biệt về cuộc thi PhysioNet/Computing in Cardiology Challenge 2013, không phải bài đề…} & {\scriptsize Bài không đề xuất mô hình nào vì đây là xã luận, và không nêu số tham số hay kích thước mô hình của bất kỳ thuật toán nào. Bài chỉ mô tả vắn tắt năm thuật toán của các đội dẫn đầu: Andreotti (ICA cộng bộ lọc…} & {\scriptsize 53 đội quốc tế tham gia, sinh ra 208 bộ nhãn và 93 bài dự thi mã nguồn mở, với điểm tốt nhất từng sự kiện là E1 = 179,439 (nhịp/phút)\textasciicircum{}2, E2 = 20,793 ms, E4 = 18,083 (nhịp/phút)\textasciicircum{}2 và E5 = 4,337 ms; Andreotti thắng E2 và E5, còn…}\newline{\scriptsize\itshape Đây là giao thức thi đấu chứ không phải kiểm chứng chéo: thí sinh huấn luyện trên Set A có nhãn công khai, ban tổ chức chấm trên Set B và Set C mà…} & một phần \\[2pt]
p21 & Adam 2020\newline{\scriptsize Đây là bài mô tả dữ liệu (Data Descriptor) trong đó nhóm Ba Lan công bố công khai hai bộ tín hiệu điện tim thai nhi do…} & {\scriptsize Bài không đề xuất mô hình học máy nào và không nêu số tham số của mô hình nào, vì đây là bài mô tả dữ liệu. Các thuật toán được nhắc đến (trừ mẫu PQRST, lọc phối hợp chuẩn hóa, Pan-Tompkins sửa đổi, PCA, ICA,…} & {\scriptsize Phần Technical Validation trích lại kết quả của chính nhóm trên bộ dữ liệu này, tất cả với dung sai ±40 ms: toàn bộ chuỗi xử lý tín hiệu bụng đạt PI tốt nhất 81,4\% nhưng khác biệt rất lớn giữa các chuyển đạo bụng (từ 70,4\% đến…}\newline{\scriptsize\itshape Bài không định nghĩa giao thức train/test nào, không tách theo chủ thể, không tách theo bản ghi và không chia ngẫu nhiên, vì đây là bài công bố dữ…} & một phần \\[2pt]
\multicolumn{5}{@{}l}{\cellcolor{rowbg}\textbf{Lý thuyết TDA}}\\[1pt]
p22 & Jose 2013\newline{\scriptsize Bài xây dựng khung lý thuyết đo tính tuần hoàn của chuỗi thời gian bằng cách kết hợp nhúng cửa sổ trượt (time-delay…} & {\scriptsize Bài không nêu số tham số vì đây không phải mô hình học máy và không có tham số huấn luyện; chỉ có siêu tham số hình học gồm số chiều nhúng M+1, bậc cắt cụt Fourier N, trường hệ số F\_p và số chu kỳ giả định L,…} & {\scriptsize Thí nghiệm 1 xếp hạng 10 tín hiệu theo điểm tuần hoàn: SW1PerS cho điểm từ 0,9488 (cosin thuần) giảm dần qua 0,8810 (ba tín hiệu), 0,8132, 0,7455, 0,6777, 0,6770 xuống 0,4066 (hai tín hiệu cuối), trong khi JTK\_CYCLE trả p-value…}\newline{\scriptsize\itshape Không có chia train/test, không có tách theo chủ thể, không có tách theo bản ghi, vì SW1PerS không học tham số nào mà chỉ tính một điểm số rồi vẽ…} & \xau{không} \\[2pt]
p23 & Henry 2017\newline{\scriptsize Sơ đồ bền vững là một đa tập điểm trong mặt phẳng với số điểm thay đổi tùy mẫu, không cộng trừ được nên không dùng…} & {\scriptsize Đây không phải mạng nơ-ron mà là một phép trích xuất đặc trưng; các bộ phân loại đặt bên trên PI gồm K-medoids với K = 6 và 1.000 lần khởi tạo ngẫu nhiên (cải thiện bằng lặp Voronoi), SVM tuyến tính chính quy…} & {\scriptsize Bảng 1 (ma trận khoảng cách 150×150, phân cụm K-medoids, độ chính xác và thời gian, mức nhiễu 0,05 và 0,1): sơ đồ gốc PD H0 với L1 đạt 96,0\% trong 37.346 s và 96,0\% trong 42.613 s, PD H0 L2 đạt 91,3\% trong 24.656 s, PD H0 L∞ đạt…}\newline{\scriptsize\itshape Đây là điểm yếu nhất của bài: thí nghiệm chính (mục 6.1, Bảng 1) dùng K-medoids, tức phân cụm không giám sát chạy trên cả 150 đám mây mà không có…} & \xau{không} \\[2pt]
p24 & Mathieu 2020\newline{\scriptsize Bài có hai đóng góp; thứ nhất là định nghĩa PersLay, một lớp mạng nơ-ron khả vi tự động nhận đầu vào là sơ đồ bền vững…} & {\scriptsize Mạng rất nhỏ, chỉ dùng 2 lớp (PersLay cộng một lớp fully-connected sau chuẩn hóa batch), và tác giả nói rõ chọn kiến trúc đơn giản để tạo ra hiểu biết chứ không nhằm đạt điểm cao nhất; cấu hình từng kênh được…} & {\scriptsize Bảng 1 (ORBIT, độ chính xác ± độ lệch chuẩn): trên ORBIT5K, Persistence Scale Space Kernel đạt 72,38 (±2,4), Persistence Weighted Gaussian Kernel 76,63 (±0,7), Sliced Wasserstein Kernel 83,6 (±0,9), Persistence Fisher Kernel…}\newline{\scriptsize\itshape ORBIT5K và ORBIT100K chia ngẫu nhiên 70\% huấn luyện và 30\% kiểm thử rồi báo cáo độ chính xác kiểm thử trung bình trên 100 lần chạy (làm giống Le và…} & \xau{không} \\[2pt]
p25 & Anirudh 2020\newline{\scriptsize Bài giải quyết nút thắt lớn nhất của TDA khi dùng trong học sâu: việc tính Persistence Diagram rồi chuyển thành…} & {\scriptsize Signal PI-Net gồm 4 lớp Conv1D với 128, 256, 512 và 1024 bộ lọc, hạt nhân 3, kèm BatchNorm, ReLU và MaxPool, kết thúc bằng global-average-pooling và lớp Dense kích thước 2500 x n rồi reshape thành 50 x 50 x…} & {\scriptsize F1 có trọng số trên GENEactiv với khung 250 bước: MLP-PI 46,45 ± 0,32; MLP-Signal PI-Net 49,47 ± 0,69; MLP-SF 41,70 ± 0,41; CNN 1D 53,56 ± 0,31; CNN 1D + PI 54,28 ± 0,23; CNN 1D + Signal PI-Net 54,41 ± 0,21; với khung 500 bước,…}\newline{\scriptsize\itshape Giao thức không đồng nhất giữa hai bộ dữ liệu và đây là điểm yếu nghiêm trọng nhất: với GENEactiv, khoảng 75\% khung được dùng để huấn luyện và phần…} & \xau{không} \\[2pt]
p26 & Renata 2023\newline{\scriptsize Đây là bài tự phản biện mạnh nhất trong cộng đồng TDA: đồng điều bền vững (PH) đã được dùng thành công ở hàng trăm ứng…} & {\scriptsize PH không phải mạng neural nên không có số tham số; bộ phân loại và hồi quy là SVM/SVR của sklearn với C thuộc \{0,001; 1; 100\}, tìm tham số bằng GridSearchCV 3 fold. Các đường cơ sở gồm SVM trên ma trận khoảng…} & {\scriptsize Đếm số lỗ hổng, độ chính xác trên dữ liệu gốc: PH 1,00; PH simple 0,94; ML 0,67; NN nông 0,53; NN sâu 0,50; PointNet 1,00 — nhưng khi tịnh tiến, PointNet sụp còn 0,20 trong khi PH giữ 1,00, khi xoay PointNet còn 0,74, và với…}\newline{\scriptsize\itshape Vì là dữ liệu tổng hợp nên không có khái niệm chủ thể hay bệnh nhân và không có nguy cơ rò rỉ dữ liệu kiểu tách theo chủ thể; họ chia ngẫu nhiên…} & \xau{không} \\[2pt]
\multicolumn{5}{@{}l}{\cellcolor{rowbg}\textbf{Nhúng trễ và chọn tham số}}\\[1pt]
p27 & Quoc 2019\newline{\scriptsize Đây là bài TDA cho chuỗi thời gian 1 chiều, gần với miền dữ liệu của nhóm nhất trong ba bài: để dùng TDA cho chuỗi…} & {\scriptsize Không phải mạng neural mà là SVM trên hạt nhân PD ba chiều, nên không có tham số học được theo nghĩa mạng, và bài không nêu số tham số cũng không nêu kích thước mô hình. Siêu tham số gồm K = 10 giá trị độ trễ…} & {\scriptsize Độ chính xác phân loại (\%) của delay-variant so với single-delay: ECG200 90,0 so với 87,0 (1NN Euclid và 1NN DTW cùng 88,0); ECG5000 93,6 so với 92,1; ECGFiveDays 99,9 so với 92,0; TwoLeadECG 99,4 so với 94,0; EigenWorms 83,1 so…}\newline{\scriptsize\itshape Với Hes1, họ chia ngẫu nhiên 50/50, lặp 100 lần chia và báo cáo độ chính xác trung bình — hợp lệ vì 1.000 tế bào là các mô phỏng độc lập; với tám bộ…} & một phần \\[2pt]
p28 & Eugene 2023\newline{\scriptsize Đây là bài tổng quan kèm đề xuất phương pháp thuộc lĩnh vực hệ động lực phi tuyến, không phải công trình y sinh: khi…} & {\scriptsize SToPS bản thân không phải mô hình học máy mà là thuật toán chọn tham số, không có tham số huấn luyện; bộ dự báo dùng để kiểm chứng là mạng truyền thẳng 4 lớp với kiến trúc 3:128:128:3, kích hoạt ReLU,…} & {\scriptsize Về phát hiện thang thời gian: với chuỗi tổng sóng sin (kỳ vọng τ = 250, 50, 8), thông tin tương hỗ chỉ tìm được thành phần tần số cao nhất, PECUZAL chỉ ra một độ trễ sai τ = 341 và MDOP ra τ = (241, 42, 271, 180) với chỉ hai giá…}\newline{\scriptsize\itshape Không có khái niệm tách theo chủ thể hay tách theo bản ghi vì không có đối tượng, chỉ có ba chuỗi tín hiệu; giao thức là dùng nửa đầu chuỗi để huấn…} & \xau{không} \\[2pt]
\multicolumn{5}{@{}l}{\cellcolor{rowbg}\textbf{TDA cho tín hiệu tim}}\\[1pt]
p29 & Yonglian 2023\newline{\scriptsize Bài giải bài toán đánh giá chất lượng tín hiệu ECG cho thiết bị đeo, tức phân loại nhị phân một đoạn ECG là "chấp nhận…} & {\scriptsize GoogLeNet (Inception v1) dùng học chuyển tiếp; bài không nêu số tham số. Bài cũng không nêu số epoch, learning rate, batch size, bộ tối ưu, cách thay lớp đầu ra, số lớp đóng băng, thư viện tính đồng điều bền…} & {\scriptsize Trên bộ 12 chuyển đạo, SubLevel-Set tốt nhất với mAcc 98,04 ± 0,11 | F1 98,40 ± 0,06 | Se 97,15 ± 0,11 | Sp 98,93 ± 0,23 (Acc 97,70), còn Vietoris-Rips giảm dần theo cửa sổ từ mAcc 95,16 ± 0,20 ở 1 giây xuống 93,11 ± 0,19; 91,75…}\newline{\scriptsize\itshape Kiểm định chéo 10 mảnh với nguyên văn "All the segments were randomly divided into 10 groups", tức chia ngẫu nhiên theo đoạn, không tách theo chủ…} & một phần \\[2pt]
p30 & Meryll 1906\newline{\scriptsize Bài giải bài toán phát hiện và phân loại loạn nhịp tim từ ECG người lớn, với trọng tâm không phải điểm số mà là khả…} & {\scriptsize Mạng nơ-ron mô-đun đa kênh gồm các kênh CNN 1 chiều và MLP nối vào một mạng kết nối đầy đủ chung, cộng một bộ tự mã hóa 6 lớp ẩn (đầu vào 400, không gian ẩn 20); bài không nêu số tham số. Bài cũng không nêu…} & {\scriptsize Mọi con số đều là độ chính xác có trọng số chứ không phải F1: phát hiện loạn nhịp nhị phân đạt 98\% trên tập kiểm định và 90\% trên bệnh nhân mới, còn phân loại 13 lớp đạt 97,3\% kiểm định và 80,5\% kiểm tra — khoảng cách 8 điểm và…}\newline{\scriptsize\itshape Điểm mạnh nhất của bài là kiểm định chéo theo bệnh nhân: tập huấn luyện, kiểm định và kiểm tra là các bệnh nhân khác nhau và được hoán vị qua toàn…} & một phần \\[2pt]
p31 & Yu-Min 1908\newline{\scriptsize Bài đặt câu hỏi liệu có thể dùng đồng điều bền vững để định lượng biến thiên nhịp tim (HRV) hay không: nhóm tác giả…} & {\scriptsize SVM nhân tuyến tính, dùng hàm fitcsvm mặc định của Matlab 2014b; khi có hơn 2 lớp thì dùng ECOC một-đối-một qua hàm fitcecoc mặc định. Đây không phải mạng nơ-ron nên không có khái niệm "số tham số huấn luyện…} & {\scriptsize Bài toán Thức vs Ngủ, huấn luyện trên CGMH-training (Bảng 1): trên CGMH-validation đạt SE 70,9 ± 16,0\%, SP 78,9 ± 5,4\%, Acc 75,8 ± 4,4\%, PR 38,1 ± 19,6\%, F1 0,452 ± 0,140, AUC 0,824 ± 0,090, Kappa 0,322 ± 0,123; trên DREAMS đạt…}\newline{\scriptsize\itshape Kiểm định chéo cơ sở dữ liệu (cross-database validation) — mạnh hơn cả tách theo chủ thể: huấn luyện trên một cơ sở dữ liệu và kiểm tra trên toàn bộ…} & một phần \\[2pt]
p32 & Andy 2025\newline{\scriptsize Đây là một nghiên cứu mô tả chứ không phải bài toán học máy: nhóm tác giả nhúng chuỗi khoảng RR của trẻ em từ 1 tháng…} & {\scriptsize Bài không có mô hình học máy nào: không có mạng nơ-ron, không có bộ phân loại, không có số tham số. Toàn bộ "mô hình" chỉ là Ripser (Python) tính sơ đồ bền vững cùng persistence landscape và các kiểm định…} & {\scriptsize Chỉ số HRV cổ điển (Bảng 2, trung vị [khoảng tứ phân vị]): Mean RR tăng đơn điệu từ 404 ms [395-417] ở sơ sinh lên 674 ms [669-680] ở thiếu niên, SDNN từ 62,9 ms [57-66] ở sơ sinh xuống thấp nhất 36,7 ms ở nhũ nhi muộn rồi lên…}\newline{\scriptsize\itshape Bài không có giao thức đánh giá học máy: không chia train/test, không tách theo chủ thể, không tách theo bản ghi, không chia ngẫu nhiên — đơn giản…} & \xau{không} \\[2pt]
p33 & Alperen 2021\newline{\scriptsize Bài xây dựng một quy trình TDA tổng quát để phân loại chuỗi thời gian đơn biến, rồi áp dụng vào bài toán nhận diện…} & {\scriptsize Bốn bộ phân loại cổ điển từ scikit-learn: Random Forest và AdaBoost Decision Tree với 100 cây và độ sâu tối đa 5; Support Vector Classifier nhân tuyến tính với tham số điều chuẩn C = 0,1; và Linear…} & {\scriptsize WESAD ba lớp (Bảng 2, tách theo chủ thể, cửa sổ 60 giây): gộp tất cả kênh với SVC và dùng tất cả sơ đồ đạt 81,35\% độ chính xác, F1 73,44\% (bản gốc WESAD 79,57\%); riêng từng kênh thì ECG 76,27\% (nhúng trễ) và 72,72\% (tất cả sơ…}\newline{\scriptsize\itshape Tách theo chủ thể (leave-one-subject-out cross-validation) cho tất cả các bộ dữ liệu — huấn luyện trên tất cả mọi người trừ một, kiểm tra trên người…} & một phần \\[2pt]
\end{longtable}